\section{Fase 1: Editor gráfico de segmentos}
\label{sec:Primera-Fase}

El objetivo de esta fase es desarrollar un editor gráfico, que nos permita crear los cartuchos con los datos del juego, de forma sencilla y visual, sin editar ficheros de texto con datos relacionados entre si, que son más difíciles de mantener.

Debemos partir de una imagen con todos los segmentos que presentaremos en pantalla para el juego, sin el fondo fijo, y en blanco y negro. El proceso descompone la imagen en los segmentos individuales, permite darles propiedades y relaciones entre ellos, unir varios en un solo gráfico, y generar el fichero JSON de datos que representa la parte gráfica del cartucho que usará el motor del juego.

Aunque el objetivo parece simple, es una fase fundamental del proyecto, ya que es importante que el usuario que quiera definir un juego lo maneje de la forma más sencilla y eficiente posible, eliminando manipulaciones manuales y ayudando en lo posible a la correcta definición de los segmentos.

